% GENERATED FILE. Edit canonical lessons, then run npm run generate:books.
% canonical-source-hash: fnv1a64:a846dd9c48fbf0ed
% canonical-lessons: MR-C551-pratiksa, MR-C551-yadi, MR-C551-kram, MR-C551-rakhiv, MR-C551-akasmik

\chapter{Waiting Lists}
\label{ch:mr-waiting-lists}
\hlchaptermodality{551}

\begin{chapteropening}
\textbf{By the end of this chapter:} \emph{I can talk about waiting, lists and turns.}

Complete the last lesson of chapter 551: I can talk about waiting, lists and turns.
\end{chapteropening}

\section[\texorpdfstring{\mr{प्रतीक्षा} (pratīkṣā)}{प्रतीक्षा (pratīkṣā)}]{\mr{प्रतीक्षा} (pratīkṣā) — waiting; expectation}

\label{lesson:MR-C551-pratiksa}



\begin{quote}
Before the new one: say the Marathi for a coin toss, then the Marathi for to bet.
\end{quote}

\subsection*{You'll want to know: \mr{प्रतीक्षा}}
\textbf{\mr{प्रतीक्षा}} — \emph{pratīkṣā} — \textquotedblleft{}waiting; expectation\textquotedblright{}.

\begin{quote}
\mr{पाहुण्यांची} \mr{आम्ही} \mr{खूप} \mr{वेळ} \mr{प्रतीक्षा} \mr{केली}.
\end{quote}

\begin{quote}
\emph{pāhuṇyāñcī āmhī khūp veḷ pratīkṣā kelī.}
\end{quote}

\begin{quote}
We waited a long time for the guests.
\end{quote}

Feminine.

The one awaited takes -\mr{ची}: \mr{पाहुण्यांची} \mr{प्रतीक्षा} \mr{करणे}.

\begin{cousinweb}
\textbf{\mr{प्रतीक्षा}} (\emph{pratīkṣā}) is a tatsama, taken from Sanskrit \emph{pratīkṣā}, \textquotedblleft{}expectation\textquotedblright{} (Monier-Williams).
\end{cousinweb}

\begin{grammarlens}[title={What you've built}]
One more word for saying what happens if.
\end{grammarlens}

\subsection*{Your turn}
\begin{itemize}
\raggedright
  \item \emph{Say it:} \emph{pratīkṣā}
  \item \emph{Say it:} \emph{pratīkṣā}, once more
  \item \emph{Say it:} \emph{pāhuṇyāñcī āmhī khūp veḷ pratīkṣā kelī.} — We waited a long time for the guests.
  \item \emph{Recall:} say the Marathi for a coin toss, then the Marathi for to bet, then say \emph{pratīkṣā} again
\end{itemize}

\begin{tcolorbox}[breakable,colback=teal!4,colframe=teal!35!black,title={Before you move on}]
What does \mr{प्रतीक्षा} mean? (\textquotedblleft{}Waiting; expectation\textquotedblright{}.) Say it once more.
\end{tcolorbox}
\section[\texorpdfstring{\mr{यादी} (yādī)}{यादी (yādī)}]{\mr{यादी} (yādī) — a list}

\label{lesson:MR-C551-yadi}



\begin{quote}
Before the new one: say the Marathi for to bet, then the Marathi for waiting.
\end{quote}

\subsection*{You'll want to know: \mr{यादी}}
\textbf{\mr{यादी}} — \emph{yādī} — \textquotedblleft{}a list\textquotedblright{}.

\begin{quote}
\mr{आईने} \mr{बाजारातून} \mr{आणायच्या} \mr{सामानाची} \mr{यादी} \mr{केली}.
\end{quote}

\begin{quote}
\emph{āīne bājārātūn āṇāycyā sāmānācī yādī kelī.}
\end{quote}

\begin{quote}
Mother made a list of the things to bring from the market.
\end{quote}

Feminine.

\begin{cousinweb}
\textbf{\mr{यादी}} (\emph{yādī}) is a loan from Persian \emph{yādī}, \textquotedblleft{}a memorandum\textquotedblright{}, built on \emph{yād}, \textquotedblleft{}memory\textquotedblright{}.
\end{cousinweb}

\begin{grammarlens}[title={What you've built}]
One more word for saying what happens if.
\end{grammarlens}

\subsection*{Your turn}
\begin{itemize}
\raggedright
  \item \emph{Say it:} \emph{yādī}
  \item \emph{Say it:} \emph{yādī}, once more
  \item \emph{Say it:} \emph{āīne bājārātūn āṇāycyā sāmānācī yādī kelī.} — Mother made a list of the things to bring from the market.
  \item \emph{Recall:} say the Marathi for to bet, then the Marathi for waiting, then say \emph{yādī} again
\end{itemize}

\begin{tcolorbox}[breakable,colback=teal!4,colframe=teal!35!black,title={Before you move on}]
What does \mr{यादी} mean? (\textquotedblleft{}A list\textquotedblright{}.) Say it once more.
\end{tcolorbox}
\section[\texorpdfstring{\mr{क्रम} (kram)}{क्रम (kram)}]{\mr{क्रम} (kram) — a sequence; an order}

\label{lesson:MR-C551-kram}



\begin{quote}
Before the new one: say the Marathi for waiting, then the Marathi for a list.
\end{quote}

\subsection*{You'll want to know: \mr{क्रम}}
\textbf{\mr{क्रम}} — \emph{kram} — \textquotedblleft{}a sequence; an order\textquotedblright{}.

\begin{quote}
\mr{शिक्षकांनी} \mr{मुलांची} \mr{नावं} \mr{वर्णमालेच्या} \mr{क्रमाने} \mr{वाचली}.
\end{quote}

\begin{quote}
\emph{śikṣakānnī mulāñcī nāvaṁ varṇamālecyā kramāne vāclī.}
\end{quote}

\begin{quote}
The teacher read out the children's names in alphabetical order.
\end{quote}

Masculine.

\begin{cousinweb}
\textbf{\mr{क्रम}} (\emph{kram}) is a tatsama, taken from Sanskrit \emph{krama}, \textquotedblleft{}a step, a sequence\textquotedblright{} (Monier-Williams).
\end{cousinweb}

\begin{grammarlens}[title={What you've built}]
One more word for saying what happens if.
\end{grammarlens}

\subsection*{Your turn}
\begin{itemize}
\raggedright
  \item \emph{Say it:} \emph{kram}
  \item \emph{Say it:} \emph{kram}, once more
  \item \emph{Say it:} \emph{śikṣakānnī mulāñcī nāvaṁ varṇamālecyā kramāne vāclī.} — The teacher read out the children's names in alphabetical order.
  \item \emph{Recall:} say the Marathi for waiting, then the Marathi for a list, then say \emph{kram} again
\end{itemize}

\begin{tcolorbox}[breakable,colback=teal!4,colframe=teal!35!black,title={Before you move on}]
What does \mr{क्रम} mean? (\textquotedblleft{}A sequence; an order\textquotedblright{}.) Say it once more.
\end{tcolorbox}
\section[\texorpdfstring{\mr{राखीव} (rākhīv)}{राखीव (rākhīv)}]{\mr{राखीव} (rākhīv) — reserved; kept aside}

\label{lesson:MR-C551-rakhiv}



\begin{quote}
Before the new one: say the Marathi for a list, then the Marathi for a sequence.
\end{quote}

\subsection*{You'll want to know: \mr{राखीव}}
\textbf{\mr{राखीव}} — \emph{rākhīv} — \textquotedblleft{}reserved; kept aside\textquotedblright{}.

\begin{quote}
\mr{ही} \mr{खुर्ची} \mr{वृद्धांसाठी} \mr{राखीव} \mr{आहे}.
\end{quote}

\begin{quote}
\emph{hī khurcī vṛddhāṁsāṭhī rākhīv āhe.}
\end{quote}

\begin{quote}
This chair is reserved for the elderly.
\end{quote}

\begin{cousinweb}
\textbf{\mr{राखीव}} (\emph{rākhīv}) is built on \mr{राखणे}, \textquotedblleft{}to keep\textquotedblright{}, inherited from Sanskrit \emph{rakṣati}, \textquotedblleft{}guards\textquotedblright{} (CDIAL 10547).
\end{cousinweb}

\begin{grammarlens}[title={What you've built}]
One more word for saying what happens if.
\end{grammarlens}

\subsection*{Your turn}
\begin{itemize}
\raggedright
  \item \emph{Say it:} \emph{rākhīv}
  \item \emph{Say it:} \emph{rākhīv}, once more
  \item \emph{Say it:} \emph{hī khurcī vṛddhāṁsāṭhī rākhīv āhe.} — This chair is reserved for the elderly.
  \item \emph{Recall:} say the Marathi for a list, then the Marathi for a sequence, then say \emph{rākhīv} again
\end{itemize}

\begin{tcolorbox}[breakable,colback=teal!4,colframe=teal!35!black,title={Before you move on}]
What does \mr{राखीव} mean? (\textquotedblleft{}Reserved; kept aside\textquotedblright{}.) Say it once more.
\end{tcolorbox}
\section[\texorpdfstring{\mr{आकस्मिक} (ākasmik)}{आकस्मिक (ākasmik)}]{\mr{आकस्मिक} (ākasmik) — unforeseen; sudden}

\label{lesson:MR-C551-akasmik}



\begin{quote}
Before the new one: say the Marathi for a sequence, then the Marathi for reserved.
\end{quote}

\subsection*{You'll want to know: \mr{आकस्मिक}}
\textbf{\mr{आकस्मिक}} — \emph{ākasmik} — \textquotedblleft{}unforeseen; sudden\textquotedblright{}.

\begin{quote}
\mr{आकस्मिक} \mr{खर्चासाठी} \mr{मी} \mr{थोडे} \mr{पैसे} \mr{बाजूला} \mr{ठेवले}.
\end{quote}

\begin{quote}
\emph{ākasmik kharcāsāṭhī mī thoḍe paise bājūlā ṭhevle.}
\end{quote}

\begin{quote}
I put a little money aside for unforeseen expenses.
\end{quote}

\begin{cousinweb}
\textbf{\mr{आकस्मिक}} (\emph{ākasmik}) is a tatsama, taken from Sanskrit \emph{ākasmika}, \textquotedblleft{}causeless, unforeseen\textquotedblright{} (Monier-Williams).
\end{cousinweb}

\begin{grammarlens}[title={What you've built}]
One more word for saying what happens if.
\end{grammarlens}

\subsection*{Your turn}
\begin{itemize}
\raggedright
  \item \emph{Say it:} \emph{ākasmik}
  \item \emph{Say it:} \emph{ākasmik}, once more
  \item \emph{Say it:} \emph{ākasmik kharcāsāṭhī mī thoḍe paise bājūlā ṭhevle.} — I put a little money aside for unforeseen expenses.
  \item \emph{Recall:} say the Marathi for a sequence, then the Marathi for reserved, then say \emph{ākasmik} again
\end{itemize}

\begin{tcolorbox}[breakable,colback=teal!4,colframe=teal!35!black,title={Before you move on}]
What does \mr{आकस्मिक} mean? (\textquotedblleft{}Unforeseen; sudden\textquotedblright{}.) Say it once more.
\end{tcolorbox}
